
\subsubsection{6.1 Summary of Findings}

The corpus-level measurements (H-E1, H-M1) demonstrate that fastText quality filtering alters demographic-occupation association structure in a monotonic manner correlated with filtering intensity at the tested scale (~50,000 documents per configuration). The entropy reduction (-22.41\%) and log-odds correlation (rho = 1.0) are both large in magnitude and statistically significant. These results are consistent with the hypothesis that fastText quality filtering acts as a demographic reweighting mechanism at the corpus level, though the mechanism has been characterized only for a single quality classifier (fastText), a single corpus (DCLM-POOL), a single demographic axis (binary gender-occupation from WinoBias), and at quick-run scale.

The model-level pilot (H-M2) did not demonstrate statistically significant graded correlation between corpus entropy and model logit margins. The negative control showed a large difference between the shuffled-demographic and unfiltered conditions, but this observation alone is insufficient to establish that corpus demographic structure propagates to model behavior in a graded manner.

H-M3 (downstream fairness benchmark comparison) was not executed, leaving the full causal chain from curation decisions to model fairness outcomes unverified.

\subsubsection{6.2 Limitations}

\textbf{L1: Proxy training for H-M2.} The model-level experiment used a proxy training setup rather than full Pythia-1B training at 100B tokens. The verification\_state.yaml describes this as a ''tiny 512-hidden proxy'' with hf\_trainer\_fallback, distinct from the intended Pythia-1B architecture (hidden\_size=2048, 16 layers). The primary gate was not satisfied, and the result cannot distinguish between insufficient model capacity/training budget and a genuine absence of corpus-to-model signal propagation.

\textbf{L2: H-M3 not executed.} The downstream fairness claim -- that matched-capability models from different curation paths produce distinguishable BBQ/WinoBias outcomes -- was not tested. The corpus-level contribution is independent of this claim, but the full PCFH hypothesis remains unverified.

\textbf{L3: Quick-run corpus scale.} H-E1 and H-M1 operated on approximately 50,000-document subsets per configuration, representing a small fraction of the full DCLM-POOL corpus. The effect magnitude (4.5x the gate threshold) makes complete reversal at full scale unlikely, but exact effect sizes should be treated as estimates. A background full-scale run was reported in progress (PID 2164630).

\textbf{L4: Single model family.} Corpus-level findings (H-E1, H-M1) are model-agnostic. The model-level pilot (H-M2) is specific to the proxy architecture used and has not been validated across model families.

\textbf{L5: Demographic lexicon scope.} The WinoBias-derived lexicon covers 20 occupations with binary-gender pronouns (male/female only), drawn from U.S. Bureau of Labor Statistics data. Non-binary gender categories, non-English documents, and non-occupational demographic dimensions (race, age, nationality) are not covered. The measured effects reflect this specific lexicon and cannot be assumed to generalize to other demographic dimensions without additional measurement.

\textbf{L6: Discrete configuration design.} The rho = 1.0 result in H-M1 is computed across five discrete configuration levels. Perfect rank correlation on five ordered points does not exclude the possibility that a finer-grained sweep would reveal deviations from monotonicity. A 20-level continuous sweep would provide stronger evidence regarding the functional form.

\textbf{L7: Unresolved data discrepancies.} The entropy values in the experimental records (h-e1/results.json) differ from the values in the paper's ground truth file (065\_ground\_truth.yaml) and corrected paper tables. The H-M2 Spearman rho differs between verification\_state.yaml (-0.2143) and the validation report (0.357). These discrepancies are noted but not fully reconciled.

\subsubsection{6.3 Broader Impact}

The corpus audit methodology (H-E1 + H-M1) provides a model-free approach for measuring demographic association structure changes introduced by quality filtering. If validated at full scale, such a tool could be applied by corpus curators before committing to model training. The methodology could also in principle be used to engineer rather than detect demographic associations; open availability of the audit approach lowers the barrier for detection of such engineering.

